\documentclass[10pt]{article}
% Shared typographic and algebraic source.
\usepackage[letterpaper,margin=0.86in]{geometry}
\usepackage[T1]{fontenc}
\usepackage[utf8]{inputenc}
\usepackage{lmodern}
\usepackage{microtype}
\usepackage{amsmath,amssymb,amsthm,mathtools,mathrsfs}
\usepackage{booktabs,array,tabularx,longtable}
\usepackage{enumitem}
\usepackage[dvipsnames]{xcolor}
\usepackage{fancyhdr,titlesec}
\usepackage[hidelinks]{hyperref}

\definecolor{ink}{HTML}{1F2A33}
\definecolor{teal}{HTML}{226A73}
\definecolor{muted}{HTML}{65747E}
\definecolor{line}{HTML}{C5D1D3}
\definecolor{wash}{HTML}{F3F7F6}

\setlength{\parindent}{0pt}
\setlength{\parskip}{0.46em}
\setlength{\emergencystretch}{2em}
\allowdisplaybreaks[2]
\renewcommand{\arraystretch}{1.16}
\setlist{nosep,leftmargin=1.5em}

\pagestyle{fancy}
\fancyhf{}
\fancyhead[L]{\footnotesize\sffamily\color{muted}COPPALINI ARCHITECTURE}
\fancyhead[R]{\footnotesize\sffamily\color{muted}PHYSICAL REALITY INSIDE CONSTRAINT}
\fancyfoot[C]{\footnotesize\sffamily\color{muted}\thepage}
\renewcommand{\headrulewidth}{0.35pt}
\renewcommand{\headrule}{\hbox to\headwidth{\color{line}\leaders\hrule height \headrulewidth\hfill}}
\setlength{\headheight}{19pt}

\titleformat{\section}{\Large\bfseries\color{ink}}{\thesection.}{0.55em}{}
\titleformat{\subsection}{\large\bfseries\color{teal}}{\thesubsection.}{0.5em}{}
\titlespacing*{\section}{0pt}{1.42em}{0.48em}
\titlespacing*{\subsection}{0pt}{1.02em}{0.32em}

\newtheorem{theorem}{Theorem}
\newtheorem{proposition}[theorem]{Proposition}
\newtheorem{lemma}[theorem]{Lemma}
\theoremstyle{definition}
\newtheorem{definition}[theorem]{Definition}
\theoremstyle{remark}
\newtheorem{remark}[theorem]{Boundary}

\newcommand{\RR}{\mathbb R}
\newcommand{\CC}{\mathbb C}
\newcommand{\QQ}{\mathbb Q}
\newcommand{\ZZ}{\mathbb Z}
\newcommand{\Id}{I}
\newcommand{\tr}{\operatorname{tr}}
\newcommand{\im}{\operatorname{im}}
\newcommand{\kerop}{\operatorname{ker}}
\newcommand{\Span}{\operatorname{span}}
\newcommand{\diag}{\operatorname{diag}}
\newcommand{\snoc}{\operatorname{snoc}}
\newcommand{\Admit}{\operatorname{Admit}}
\newcommand{\Gate}{\operatorname{Gate}}
\newcommand{\Typed}{\operatorname{Typed}}
\newcommand{\Future}{\operatorname{future}}
\newcommand{\Scope}{\operatorname{sc}}
\newcommand{\SilScope}{\mathsf{Silence}_{\rm scope}}
\newcommand{\SilFit}{\mathsf{Silence}_{\rm fit}}

\newcommand{\papertitle}[2]{%
  \begin{center}
  {\color{teal}\rule{\textwidth}{1.35pt}}\\[0.38in]
  {\footnotesize\sffamily\bfseries\color{teal} COPPALINI ARCHITECTURE}\\[0.16in]
  {\fontsize{31}{36}\selectfont\bfseries\color{ink}#1\par}
  \vspace{0.15in}
  {\Large\color{teal}#2\par}
  \vspace{0.22in}
  {\color{line}\rule{\textwidth}{0.55pt}}\\[0.12in]
  {\normalsize Anthony Vito Coppa}\par
  {\small 12 November 2026}\par
  \vspace{0.30in}
  \end{center}
}

\newcommand{\sharedroot}{%
Let
\[
Q=\begin{pmatrix}1&1\\1&0\end{pmatrix},\qquad
\Omega=\begin{pmatrix}0&1\\-1&0\end{pmatrix},\qquad
C=-\Omega,
\]
\[
G=[Q,\Omega]=\begin{pmatrix}-2&1\\1&2\end{pmatrix},\qquad
H=\frac{G}{\sqrt5},\qquad
S=\frac{2Q-I}{\sqrt5}.
\]
The common exact identities are
\[
Q^2=Q+I,\qquad C^2=-I,\qquad G^2=5I,\qquad H^2=S^2=I,
\]
\[
H=CS,\qquad SG=\sqrt5\Omega=-GS,\qquad
HC=-CH,\qquad HCH=-C,
\]
\[
L_s:=e^{sS},\qquad
L_s^{\mathsf T}GL_s=G,\qquad HL_sH=L_{-s}.
\]
}

\newcommand{\sharedoccurrence}[2]{%
A physical field has an interface
\[
\begin{aligned}
\mathfrak I_J:=(&J,\Gamma_J,\operatorname{Src}_J;
A_J,X_J,Y_J,T_J,\Sigma_J,\operatorname{Con}_J,\\
&\rho_J,O_J,K_J,B_J,\theta_J),
\qquad
\rho_J:X_J\to Y_J,\quad
O_J:X_J\to\mathcal O_J,\quad
\theta_J:X_J\to\{0,1\}.
\end{aligned}
\tag{#1}
\]
Its occurrence type is
\[
\begin{aligned}
\mathsf{Occ}_J:=\{&
(\mathfrak I_J;\Theta,x,\rho_J(x),\kappa,
\boldsymbol\varepsilon,\theta_J(x)):\\
&\Theta\in A_J,\ x\in X_J,\ \kappa\in K_J,\
\boldsymbol\varepsilon\in\RR_{\ge0}^{3}\}.
\end{aligned}
\tag{#2}
\]
An indexed occurrence $\mathsf{Occ}_{J,n}$ is an element of
$\mathsf{Occ}_J$.
}

\newcommand{\sharedmatching}{%
Let
\[
\mathcal A_8:=\{\infty\}\sqcup\ZZ_7,
\qquad
(\infty,0,1,2,3,4,5,6)
=(\mathsf{Gr},\mathsf{Li},\mathsf{Me},\mathsf{Ti},
\mathsf{Hz},\mathsf{Si},\mathsf{Sc},\mathsf{Mo}).
\]
For $r\in\ZZ_7$, set
\[
\mathsf M_r=\{\{\infty,r\}\}\cup
\{\{r+j,r-j\}:j=1,2,3\}.
\]
Then $E(K_8)=\bigsqcup_{r\in\ZZ_7}\mathsf M_r$, and the canonical matching is
\[
\mathsf M_0=\{\mathsf{Gravity\! -\! Light},\;
\mathsf{Memory\! -\! Motion},\;
\mathsf{Time\! -\! Scope},\;
\mathsf{Horizon\! -\! Silence}\}.
\]
}

\hypersetup{pdftitle={Suture Field: Solder Flux and Metric g},pdfauthor={Anthony Vito Coppa},pdfsubject={Eight-address carrier, Lorentzian suture, induced flux, and Einstein-Maxwell realization}}
\newcommand{\hateta}{\widehat\eta}
\newcommand{\hatC}{\widehat C}
\newcommand{\hatrho}{\widehat\varrho}
\newcommand{\hatH}{\widehat H}
\newcommand{\Fsol}{\mathcal F_{\rm sol}}
\newcommand{\Ubar}{\overline{\mathsf U}}
\begin{document}
\papertitle{Suture Field}{Solder flux and metric $g$}

\begin{abstract}
Spacetime interval and oriented surface flux arise from one rank-four field.
The field sutures the tangent carrier of a Lorentzian manifold into an
eight-address internal carrier. Its symmetric contraction is the metric $g$;
its antisymmetric contraction is the solder flux $\Fsol$. The sevenfold
address action preserves both contractions, while the mirror preserves the
metric and reverses flux orientation. Einstein curvature and Maxwell curvature
meet when the induced flux satisfies the Maxwell field equations under a
declared dimensional calibration.
\end{abstract}

\section{Suture relation}

Let $M$ be a smooth four-manifold. Let $E\to M$ be a real rank-eight vector
bundle whose fibers carry a symmetric form $\hateta$, an antisymmetric form
$\hatC$, a cyclic action $\hatrho(\ZZ_7)$, and a mirror $\hatH$. A suture
field is a smooth $E$-valued one-form
\[
\Pi=\Pi^A{}_{\mu}\,dx^\mu\otimes e_A\in\Omega^1(M,E)
\]
of constant rank four such that
\[
g_{\mu\nu}:=\hateta_{AB}\Pi^A{}_{\mu}\Pi^B{}_{\nu}
\quad\text{has signature }(3,1).
\tag{SF.0}
\]
Equivalently, at every address $x\in M$,
\[
\Pi_x:T_xM\hookrightarrow E_x,
\qquad
g_x(u,v)=\hateta_x(\Pi_xu,\Pi_xv).
\tag{SF.1}
\]
Put $L_\Pi:=\im\Pi\subset E$. The rank and signature conditions make $L_\Pi$ a
Lorentzian rank-four subbundle and give a fiberwise isometry
\[
\Pi:TM\xrightarrow{\ \cong\ }L_\Pi.
\tag{SF.2}
\]
Thus $\Pi$ is an extended solder coframe. The word \emph{suture} names the
same coframe together with its four internal transverse address directions.

\section{Eight-address carrier}

\sharedmatching

Let
\[
\mathcal V_8:=\RR[\mathcal A_8],
\qquad
\tau e_\infty=e_\infty,\quad \tau e_k=e_{k+1}.
\tag{SF.3}
\]
The shift permutes the matchings: $\tau\mathsf M_r=\mathsf M_{r+1}$.
The cyclic action is read in the real Fourier decomposition of
$\mathcal V_8$. Let
\[
u_0=\frac1{\sqrt7}\sum_{k=0}^{6}e_k,
\]
and, for $j=1,2,3$, let
\[
\begin{aligned}
c_j&=\sqrt{\frac27}\sum_{k=0}^{6}
\cos\!\left(\frac{2\pi jk}{7}\right)e_k,\\
s_j&=\sqrt{\frac27}\sum_{k=0}^{6}
\sin\!\left(\frac{2\pi jk}{7}\right)e_k.
\end{aligned}
\tag{SF.4}
\]
Then
\[
\mathcal V_8=\Span(e_\infty,u_0)\oplus
\bigoplus_{j=1}^{3}\Span(c_j,s_j).
\tag{SF.5}
\]
Writing $U_\alpha=e^{\alpha C}$, the address action becomes
\[
\hatrho(\tau^r)=I_2\oplus U_{2\pi r/7}\oplus
U_{4\pi r/7}\oplus U_{6\pi r/7},
\qquad r\in\ZZ_7.
\tag{SF.6}
\]
The first $I_2$ acts on $\Span(e_\infty,u_0)$; the three rotations act on
the ordered pairs $(c_j,s_j)$. Hence SF.6 is a real eight-dimensional
realization of the vertex permutation action in SF.3.

\sharedroot

On the Fourier carrier, define
\[
\hateta:=G\oplus I_2\oplus I_2\oplus I_2,
\qquad
\hatC:=C\oplus C\oplus C\oplus C,
\tag{SF.7}
\]
and, with $D=\diag(1,-1)$,
\[
\hatH:=H\oplus D\oplus D\oplus D.
\tag{SF.8}
\]
The symmetric form has signature $(7,1)$. The first summand
$\Span(e_\infty,u_0)$ is the $\mathsf{Gr}$--$\Ubar$ carrier: gravity
$e_\infty=\mathsf{Gr}$ and the cycle mean $\Ubar:=u_0$. The address
$e_0=\mathsf{Li}$ has weight $1/7$ on $u_0$ and $6/7$ on the transverse planes,
so $\mathsf{Li}$ is not a summand of its own. The remaining three planes carry
the nontrivial real Fourier weights of the sevenfold address cycle.

\begin{theorem}[Cyclic and mirror invariance]
For every $r\in\ZZ_7$,
\[
\hatrho(\tau)^7=I_8,
\qquad
\hatrho(\tau^r)^{\mathsf T}\hateta\hatrho(\tau^r)=\hateta,
\qquad
\hatrho(\tau^r)^{\mathsf T}\hatC\hatrho(\tau^r)=\hatC,
\tag{SF.9}
\]
and
\[
\hatH^{\mathsf T}\hateta\hatH=\hateta,
\qquad
\hatH^{\mathsf T}\hatC\hatH=-\hatC,
\qquad
\hatH\hatrho(\tau^r)\hatH=\hatrho(\tau^{-r}).
\tag{SF.10}
\]
\end{theorem}

\begin{proof}
The three transverse blocks have eigenvalues
$e^{\pm2\pi i r/7}$, $e^{\pm4\pi i r/7}$, and
$e^{\pm6\pi i r/7}$, while the root block is fixed. Thus seven cyclic
steps give $I_8$. Every $U_\alpha$ preserves both $I_2$ and $C$.
The shared root gives $H^{\mathsf T}GH=G$ and $HCH=-C$; each $D$ gives
$D^{\mathsf T}I_2D=I_2$ and $DCD=-C$. Finally, $DU_\alpha D=U_{-\alpha}$.
Taking direct sums gives SF.9--SF.10.
\end{proof}

\section{Metric $g$}

For a locally constant $r:U\to\ZZ_7$, the internal cyclic update is
\[
\Pi\longmapsto\Pi^{(r)}:=\hatrho(\tau^{r})\Pi.
\tag{SF.11}
\]
The metric is invariant:
\[
\begin{aligned}
g^{(r)}_{\mu\nu}
&=\hateta_{AB}(\Pi^{(r)})^A{}_{\mu}(\Pi^{(r)})^B{}_{\nu}\\
&=\hateta_{AB}\Pi^A{}_{\mu}\Pi^B{}_{\nu}
=g_{\mu\nu}.
\end{aligned}
\tag{SF.12}
\]
The mirror does not commute with the cyclic action. By SF.10,
$\hatH\hatrho(\tau^w)=\hatrho(\tau^{-w})\hatH$, so $\hatH$ intertwines $\hatrho$
with $\hatrho\circ\sigma$. In the gauged case of \S6 (principal $\ZZ_7$-bundle
$P\to M$, $E=P\times_{\hatrho}\mathcal V_8$), set
\[
\sigma(\tau^w)=\tau^{-w},\qquad
E^\sigma:=P\times_{\hatrho\circ\sigma}\mathcal V_8,\qquad
\bar\Pi:=\hatH\Pi\in\Omega^1(M,E^\sigma),\qquad
\bar\Pi^{(w)}=\hatrho(\tau^{-w})\bar\Pi=\hatH\Pi^{(w)}.
\tag{SF.13}
\]
Both bundles carry the same $\hateta$ and $\hatC$, and
$\hatH:(E,\hateta,\hatC)\to(E^\sigma,\hateta,-\hatC)$ is an isometry that
reverses $\hatC$. In particular
\[
g(\bar\Pi)=g(\Pi).
\tag{SF.13a}
\]
The Levi--Civita connection and curvature are therefore determined by the
suture metric:
\[
\Gamma^\lambda{}_{\mu\nu}
=\frac12g^{\lambda\sigma}
(\partial_\mu g_{\sigma\nu}+\partial_\nu g_{\sigma\mu}
-\partial_\sigma g_{\mu\nu}),
\qquad
\operatorname{Ein}_{\mu\nu}
=R_{\mu\nu}-\frac12Rg_{\mu\nu}.
\tag{SF.14}
\]

\section{Solder flux}

The antisymmetric contraction of the same field is the solder flux
\[
\Fsol:=\frac12\hatC_{AB}\Pi^A\wedge\Pi^B
\in\Omega^2(M).
\tag{SF.15}
\]
For every oriented surface $\Sigma\subset M$, its integrated flux is
\[
\Phi_{\rm sol}(\Sigma):=\int_\Sigma\Fsol.
\tag{SF.16}
\]
The same cyclic update that preserves $g$ preserves $\Fsol$, while the
mirror reverses its orientation:
\[
\Fsol(\Pi^{(r)})=\Fsol(\Pi),
\qquad
\Fsol(\bar\Pi)=-\Fsol(\Pi).
\tag{SF.17}
\]
Metric and flux thus arise at the same address from the symmetric and
antisymmetric contractions of $\Pi$.

\section{Electromagnetic and gravitational realization}

Let $\chi$ be a declared constant dimensional calibration. An electromagnetic
realization of solder flux is the field
\[
F_{\rm em}:=\chi\Fsol
\tag{SF.18}
\]
when it satisfies
\[
dF_{\rm em}=0,
\qquad
d\star_gF_{\rm em}=\mu_0\star_gJ_{\rm em}.
\tag{SF.19}
\]
On a contractible region, the first condition gives a potential
$A$ with $F_{\rm em}=dA$. Its $U(1)$ loop read is
\[
W^{U(1)}_\gamma=
\exp\!\left(\frac{iq}{\hbar}\oint_\gamma A\right)
=\exp\!\left(\frac{iq}{\hbar}\int_\Sigma F_{\rm em}\right)
\quad(\partial\Sigma=\gamma).
\tag{SF.20}
\]
The electromagnetic stress--energy tensor is
\[
T^{\rm em}_{\mu\nu}=\frac1{\mu_0}
\left(F_{{\rm em}\,\mu\alpha}F_{\rm em\,\nu}{}^\alpha
-\frac14g_{\mu\nu}F_{{\rm em}\,\alpha\beta}
F_{\rm em}^{\alpha\beta}\right).
\tag{SF.21}
\]
With material stress--energy $T^{\rm mat}_{\mu\nu}$, curvature obeys
\[
\operatorname{Ein}_{\mu\nu}+\Lambda g_{\mu\nu}
=\frac{8\pi G_N}{c^4}
\left(T^{\rm mat}_{\mu\nu}+T^{\rm em}_{\mu\nu}\right).
\tag{SF.22}
\]
Thus the suture metric supplies interval and curvature, and the calibrated
solder flux supplies the electromagnetic two-form and its energy--momentum.

\section{Sevenfold holonomy}

Let $P\to M$ be a principal $\ZZ_7$ bundle with associated internal bundle
$E=P\times_{\hatrho}\mathcal V_8$. If a closed loop $\gamma$ has monodromy
$\tau^{w_\gamma}$, $w_\gamma\in\ZZ_7$, its internal holonomy is
\[
W^{(7)}_\gamma=\hatrho(\tau^{w_\gamma}).
\tag{SF.23}
\]
The mirror field $\bar\Pi$ of SF.13 is a section of $E^\sigma$; its holonomy
around $\gamma$ is $\hatrho(\tau^{-w_\gamma})$.
The three nontrivial Fourier planes read its phase as
\[
\exp\!\left(\frac{2\pi i j w_\gamma}{7}\right),
\qquad j=1,2,3.
\tag{SF.24}
\]
The $\ZZ_7$ holonomy and the $U(1)$ holonomy have distinct carriers. They
meet in one physical realization when a material coupling and a calibrated
observable relate $W^{(7)}_\gamma$ to $F_{\rm em}$ at the same address.

\section{Direct occurrence}

\sharedoccurrence{SF.25}{SF.26}

For $J=\mathsf{SF}$, take the state coordinates to include
\[
x_n=(\Pi_n,g_n,(\Fsol)_n,F_{{\rm em},n},J_{{\rm em},n},W^{(7)}_{\gamma,n}),
\tag{SF.27}
\]
whenever the Maxwell conditions in SF.19 have been supplied. The suture
projection is
\[
\operatorname{pr}_{\rm SF}(\mathsf{Occ}_{\mathsf{SF},n})
=(\Theta_n,g_n,(\Fsol)_n,F_{{\rm em},n},J_{{\rm em},n},
W^{(7)}_{\gamma,n},O_n,\kappa_n,\boldsymbol\varepsilon_n).
\tag{SF.28}
\]
The metric, flux, calibration, current, loop read, observable, and residual
remain separately measured coordinates of one occurrence.

\section{Suture closure}

At every spacetime address, the field closure is
\[
\boxed{
\Pi\in\Omega^1(M,E)
\quad\longmapsto\quad
\left(
g=\Pi^*\hateta,\;
\Fsol=\Pi^*\hatC
\right).}
\tag{SF.29}
\]
Here the first pullback is symmetric and the second is antisymmetric:
\[
\begin{array}{rcl}
\text{Suture Field }\Pi&:&TM\xrightarrow{\ \cong\ }L_\Pi\subset E,\\[2pt]
\text{Metric }g&:&\text{interval, connection, and curvature},\\[2pt]
\text{Solder Flux }\Fsol&:&\text{oriented surface relation and loop phase}.
\end{array}
\tag{SF.30}
\]
The cyclic address action preserves both $g$ and $\Fsol$; the mirror preserves
$g$ and reverses $\Fsol$. The Maxwell conditions select the electromagnetic
realization of flux, and the Einstein equation couples its stress--energy to
spacetime curvature.

\end{document}
